\section{15" Log}
\subsection{10/04/2023}
Juan, Hunter, and Preston began to build the wheel assemblies.

\subsection{10/13/2023}
Juan and Hunter began to assemble the chassis and the drivetrain. They then installed
the wheel assemblies, and completed the drivetrain. This was done on only one side
in order to find the ideal spacing. After the appropriate spacing was agreed upon,
they squared the chassis off and installed a few motors. At this time it is important to note that they
were missing the 3D prints that are located in the 15" cad and were using last year's motor stack
support prints which were found to work out better than what was in the cad.

\subsection{10/19/2023}
Juan and Hunter replicated the drivetrain onto the other side of the chassis. They then began to
look and find the ideal spacing so that both sides would have the least amount of friction. Eventually,
the chassis was completed, and the chassis was squared off. At this time the robot is still missing the
intended 3D prints located on the cad. They then installed the motor covers for future reference since
they were still missing two blue motors. After the completion of the chassis, Juan began to cad
the wings and the intake.
\insertimage{images/october/completed15chassis.jpg}{0.1}

\subsection{10/28/2023}
Before working on the intake, Hunter and Juan installed the front and rear chassis support blocks
to further stabilize the chassis. After this, they began to work on the first iteration of the wing assembly.
They aimed at using the cad that was created, however, it was soon found out that they would need to deviate
from the cad and change to a 3 by c channel to house the piston. This was done on only one side and is still
in the process of being finalized. They then proceeded to move on to building the intake. The first design
of the intake was directly drove with a 200rpm green motor. While building the intake, it was decided that the
cad was not going to work as expected and changes were made. They then proceeded to work on designing the
intake retraction mechanism. The first iteration came with using a 200rpm motor, but this did not provide
the necessary torque to retract, thus Juan began to cad iteration 2.
\insertimage{images/october/first iteration of intake.jpg}{0.1}

\subsection{11/01/2023}
Juan, Hunter, and Collin installed the 3D printed intake brace.

\subsection{11/02/2023}
Juan and Hunter built the second iteration of the intake retraction arm. They also began tuning the intake
and retraction arm extension. Through the tuning of the intake, it was decided to add a blue motor in place of
the green motor so that the triballs would intake more efficiently.

\subsection{11/09/2023}
After long consideration, it was decided the retraction mechanism was taking up too much space. Hunter and
Juan decided that iteration 2 needed to be slid back on the robot to allow for more space of future subsystems.
Everything else regarding the intake and the way it retracted remained the same. Also, the team decided with
this new allocated space a flywheel would be able to fit inside the robot and be the way the 15" manipulates
the triball.

\insertimage{images/november/final iteration of retraction mech.jpg}{0.1}

\subsection{11/10/2023}
Juan slid back the other side of the retraction mechanism and began to prototype a 3000rpm flex wheel flywheel
and how it could be incorporated into the robot. After seeing that the design had potential, Juan began to
cad the flywheel.

\subsection{11/16/2023}
Juan and Hunter installed new 3D printed intake supports blocks for the rear chassis that
provided both aesthetic and structure for the robot. They also noticed that the retraction arm motors
needed to be flipped in order to provide enough space for the flywheel assembly. The flywheel assembly was
then installed. At this time, another motor is still needed to power the flywheel, and they plan on switching
to regular high strength bearing flats instead of ball bearings.

\insertimage{images/november/top view of flywheel.jpg}{0.1}
\insertimage{images/november/sideview with intake extension.jpg}{0.1}

\subsubsection{11/23/2023 - 1/18/2024}
The team did not hold meetings or work on the robot due to the holidays.

\subsection{1/19/2024}
Upon returning from the break, Juan and Hunter installed the second motor for the flywheel and began testing to 
determine what was working and what was not. 

\subsection{1/21/2024}
Hunter and Juan mounted the brain and the battery and began to start routing wires on the robot. This was done 
to make flywheel testing easier. They experimented with what would have the best compression and allow the triballs 
to shoot consistently. 

\subsection{1/25/2024}
Juan, Hunter, and Braden continued flywheel testing. Braden programmed the flywheel toggle onto the controller for easier 
testing and began to play around with various rpms. After testing, they noticed that the flywheel's movement 
made the robot unstable, so they added additional rear bracing.

\insertimage{images/January/flywheel bracing.jpg}{0.1}

\subsection{1/26/2024}
The team was able to order some more blue cartridges so Hunter and Juan added the 5th and 6th drive motors which then
allowed them to incorporate the motor stacks from last year's robot used at Worlds.

\insertimage{images/January/boxed motor stack.jpg}{0.3}

\subsection{1/27/2024}
Juan began to practice ramming in triballs and the team realized that the chassis was beginning to cave in. Thus, they
decided to brace the chassis with standoff that run underneath the robot. Hunter and Juan also realized that the CAD wings
were not going to be able to work due to the 15" size constraint. They also continued tuning the flywheel.

\subsection{1/28/2024}
Seeing that the original wing design was not going to be feasible, Juan and Hunter began to prototype drop down wings
since the only way to make the cad design work would require a complete rebuild of the chassis. They decided that the
wings would be powered by two double-acting pistons and started looking at how they could be oriented on the robot.

\subsection{2/1/2024}
After prototyping, Juan went back to the cad and began looking at what might be the most reasonable way to mount the wings.
He was able to discover a more robust mounting system that would keep the robot in size and provide the wings with structural
integrity. The two of them then installed the new wing mechanism onto the robot.

\insertimage{images/february/first iteration of wings.jpg}{0.1}

\subsection{2/2/2024}
Since the new wing design looked promising, Juan mounted the air tanks on the robot so that they could be tested properly.
While testing, Juan noticed that the wing's integrity could be compromised if they were hit head on, so he braced them 
with a standoff collar joint mechanism.

\insertimage{images/february/braced wing and air tank.jpg}{0.1}

\subsection{2/3/2024}
While testing the wings, Juan and Hunter noticed that the pistons connected to the c-channel were beginning to bend.
They deduced this was due to the force of the stroke being applied when the pistons were activated. They decided to swap
out the SMC branded pistons for HLQD branded pistons since the internal piston shaft was stronger, and offered a better way 
to orient the pistons. However, the wings were still not deploying properly, so they added springs and rubber bands to help
the pistons have a more consistent stroke. The original wing design system was then removed since it was no longer needed.
Hunter and Juan also added additional support to reduce chassis torsion. The standoff bracing that was added also acted as a 
hard-stop for the intake retraction arms. This made the intake arms spin only in a semicircle and force it to stop at the
same spot every time.

\subsection{2/6/2024}
The team builders dyed the green intake sprockets black to add more of an aesthetic feature to the robots. This process
was rather simple and was done by boiling water, acetone, and black dye. The sprockets were then placed into the mixture
for a couple of minutes until the black dye had completely soaked into the gears.

\insertimage{images/february/real deal black sprockets.jpg}{0.1}



\subsection{2/7/2024}
Hunter and Juan replaced the green intake sprockets with the ones that the team dyed black. They also replaced the gears
that powered the intake arms with a 3d printed cam shaped PLA piece that mimicked the geometry of the 60 tooth gear that
was originally there. In addition, they replaced the red 3d printed intake brace with one that was black to make it look
more appealing. The team also began to mess around with the LED RGB strips and got them to work through the Vex brain. 
Juan also installed a 3d printed plate piece that covered the flywheel motors to make the robot look even nicer. 

\insertimage{images/february/rgb strips working.jpg}{0.1}
\insertimage{images/february/side view with plate.jpg}{0.1}

\subsection{2/8/2024}
Hunter and Juan mounted the license plates, wrapped the air tanks and pistons in a carbon fiber wrap, and the barrier ramps
that were in the CAD file were mounted. However, the compact nature of the 15" robot  affected the match loading capability 
of the robot. The 1/4" difference of where these ramps were placed drastically reduced the efficiency of the flywheel.
Luckily, the chassis was designed to be symmetrical, and the ramps were moved to the back of the robot. Unfortunately, the
robot was now out of size. So Juan and Hunter moved the intake hard stop had to be moved further in, so the air tanks could 
also be moved in as well. They also begin looking into potential IMU mounting positions. 

\subsection{2/9/2024}
In order to give a robot a unique finish, the team decided to anodize easily removable hardware or subsystems of the robot.
Most of these parts needed to be rebuilt into more competition ready variants, so they were going to be replaced anyway. 
Juan and Hunter disassembled these components and laid them out for an easier rebuild. The following parts were removed:

\begin{itemize}
    \item intake mechanism
    \item wings
    \item outer chassis components
    \item intake retraction bearings
    \item wing piston plate
    \item license plate mounting hardware
\end{itemize}

These parts were then cleaned thoroughly and went under a multi-stage anodization process. The removed parts
were scrubbed with dish soap, sprayed and scrubbed with a degreaser, and then rinsed off with distilled water.

\insertimage{images/february/removed c channels for anodize.jpg}{0.1}
\insertimage{images/february/15 anodize intake layout.jpg}{0.1}

\subsection{2/10/2024}
The next day, Juan, Hunter, Braden, and Collin started the next steps of the anodization process. 
Since the robot was taken apart, Juan and Hunter installed the IMUs in a suitable location
and mounted the GPS sensor. While the anodization process took many hours and the rebuild seemed to take even longer,
the team was happy with the results and how the robots turned out. 

\insertimage{images/february/anodize first step.jpg}{0.1}
\insertimage{images/february/anodize last step.jpg}{0.1}
\insertimage{images/february/anodized complete 15.jpg}{0.1}

\subsection{2/11/2024}
On this day, Hunter and Juan cnc'd the skirts out of ABS and mounted them and installed the odometers. Juan designed
an additional piece that was to be added to the wings so that they could push triballs more effectively. The team
was then donated a sheet of Lexan. Thus, the Juan and Hunter decided to cnc the pieces for the wings out of Lexan 
rather than ABS to provide more durability over time.

\insertimage{images/february/cnc skirts.jpg}{0.1}

\subsection{2/12/2024}
Since Braden wanted to create a macro that could be used for match loading, Hunter mounted the distance sensors inside the
intake of the robot. This would allow the robot to detect when triballs were placed in front of it and would know when it
could extend and retract with the press of a button. Juan then mounted a port expander at the back of the robot and cnc'd 
the aluminum ATUM plate. While this was being done, the team decided to run a couple of skills matches before leaving
the lab. Due to an extensive run time, the right intake bar snapped where it is connected to the retraction linkage.

\subsection{2/13/2024}
Everett replaced the snapped intake bar and then later that evening Juan and Hunter began working on more flywheel tuning.
While doing this, they realized that the flywheel needed to be spaced further apart to account for intake error. Rebuilding 
the flywheel did help keep the triballs straight but, in turn, reduced the compression and made the triballs not be able
to shoot over the barrier. They also reprinted the plate piece and added a piece of Lexan on top of it so the RGB lighting
could shine through the aluminum ATUM plate.

\insertimage{images/february/real deal atum.jpg}{0.1}

\subsection{2/14/2024}
Juan, Hunter, Preston and Everett worked on tuning the intake for skills with Braden and Darius and also practiced
match loading. While testing, the leftside intake bar snapped and the builders made the decision to replace both sides 
that snapped with c-channels.

\insertimage{images/february/snapped intake arm.jpg}{0.1}

\subsection{2/16/2024}
Braden and Chris continued tuning and noticed that the flywheel was not spinning properly. Upon inspection, they
found that the versa hubs were worn out due to friction caused by a cracked shaft collar. The versa hub were then
reprinted with ABS instead of PLA and were then replaced. After replacing them, Juan and Hunter continued to tune the
flywheel with Braden.

\subsection{2/17/2024}
Juan and Hunter helped Braden and Darius tune the flywheel for a few hours. This included trying out varying rpms
with different levels of intake extension and compression.


